% -----------------------------------------------------------------------------
% Conclusion
% -----------------------------------------------------------------------------
\section{Conclusion}
\label{sec:conclusion}
This study developed and evaluated a chronological sequence--ensemble
architecture for insider threat detection on the synthetic CERT benchmarks.
The leakage-checked representation kept inactive days and preserved
chronological partition boundaries, and the temporal experiments confirmed
that both the order of activity and the accumulated behavioural history
contributed to prediction.

\smallskip
Teacher-anchored logit-and-route regularisation gave stable end-to-end
optimisation of the bidirectional long short-term memory
(Bi-LSTM)--attention--oblivious differentiable sparsemax tree (ODST) student
across all three CERT r5.2 seeds. The students closely matched their frozen
teachers, supporting reproducibility, and the teacher was removed at
inference.

\smallskip
Under matched input content, ODST and attention--linear both reached a mean
PR-AUC of 0.931. Attention--linear was the stronger efficiency and
raw-calibration reference, while ODST added faithful reference-centred
tree-contribution evidence with traceable routes and leaves. Random Forest
and Extreme Gradient Boosting remained strong engineered-tabular references.
The explanation faithfulness transferred across CERT versions, grouped
temperature scaling improved probability calibration, and component
profiling identified the ODST head as the main target for runtime
optimisation.

\smallskip
Future work will extend the evaluation to CERT r6.2 once its metadata,
ground truth, feature interface, and temporal partition have been audited
and frozen, and will examine causal sequence encoders, source- and
noise-aware training, runtime optimisation, and human evaluation of the
explanation evidence.

\section*{Data Availability}
The CERT r4.2 and r5.2 insider-threat datasets used in this study are publicly
available from the Carnegie Mellon University CERT Division insider-threat
test dataset~\cite{GlasserLindauer2013} and are not redistributed here. The
derived artifacts required to reproduce the reported results---the
sequence-construction configuration, dataset manifests, and the SHA-256 hashes
of the frozen tensors---are provided in the code repository listed below; the
raw CERT data must be obtained separately from the original source under its
own licence.

\section*{Code Availability}
The source code for data preparation, teacher-anchored training, evaluation,
and the explanation analyses is available at
\url{https://github.com/alriyami7777-alt/teacher-anchored-insider-threat-detection}.
The repository includes the fixed seeds, validation-selected thresholds, and
manifests needed to reproduce the results reported in this paper.